\section{Lesson 4, 30/09}
The generators are useful because one can study the entire group just by considering
the generators: they are operators uniquely defined by their commutation relations
using Lie brackets; the set of generators equipped with their commutation relations
is the \textbf{Lie algebra} associated with the group.

Once the algebra is known, the choice of the representation yields matrices equivalent
to the generators.

As first examples, it is useful to consider $SO(3)$ and $SU(2)$, because they are
closely related to the Lorentz group (which is the final purpose of this parenthesis
on group theory).
\subsection{Example: angular momentum as a representation}
As is known, the algebra of the angular momentum operators $J^i$ is
\begin{gather*}
    \left[J^i,J^j\right] = i\varepsilon\indices{^i^j_k}J^k
\end{gather*}
the rank of a group is the maximum number of generators that commute with one another:
in this case the rank is $1$ and the chosen operator is $J^3$.\\
A Casimir operator is an operator that commutes with all generators: in this case it
is
$$J^2 = J_iJ^i.$$
The physical question is then "how many possible values/eigenstate of the angular
momentum are possible?" and it turns out that the answer to this question also
answers, in a QFT formalism, "how many particles does the field predict?".\\
The process to find the possible eigenvalues and eigenstates is known (see QM1 course)
and it yields the fact that for each $j$ value, the number of eigenstates of $J^3$
is finite and their eigenvalues $m$ must be in the range
$$-j< m < j$$
where $j$ is called the \textit{maximum weight}; so for each $j$ the possible values
of $J^3$ (and so the number of different eigenstates) are $2j+1$. From this process
also follows that the only possible values of $j$ are positive integers or
semi-integers and that the values of $m$ always differ by an integer (so if $j$ is
semi-integer, all the corresponding values of $m$ must also be semi-integer and
vice-versa).

This formalism is very powerful because it relates the irreducible representations
of the algebra to its vector basis: in a QFT, the dimension of the vector space
will correspond to the number of possible physical particles.

This algebra is common between $SU(2)$ and $SO(3)$, with the key difference that
for $SO(3)$ only integer values of $j$ are permitted: as seen in intro groups, $SU(2)$
can be represented as a 3-sphere immersed in $\mathbb{R}^4$, while $SO(3)$ is
only half 3-sphere in which antipodal points are mapped to the equator (??).

\subsection{Example: reducible representation}
The combination of 2 spin $\frac{1}{2}$ particles yields a reducible representation:
the vector space is
$$V = \mathbb{C}^2\otimes\mathbb{C}^2,\quad \dim(V)=4$$
again, the process of spin combination is treated in the QM1 and QM2 courses and it
yields a \textit{triplet} and a \textit{singlet} of states:
\begin{align*}
&\text{Triplet (J=1): }\begin{dcases*}
\ket{J=1,M=-1}\\
\ket{J=1, M=0}\\
\ket{J=1, M=-1}
\end{dcases*}\rightarrow\text{Vector states}\\
\phantom{x}\\
&\text{Singlet J=0}: \ket{J=0,M=0} \rightarrow\text{Scalar state}
\end{align*}
So the space $V$ is subdivided into two different vector spaces which differ in
symmetry: the singlet is antisymmetric while the triplet is symmetric; this
representation is reducible, because the image of the representation of the algebra
is the direct sum of the images of the representations of two different algebras.
\subsection{The Lorentz Group Representation}
Only the subgroup $L^\uparrow_+$ is here considered and its formal name is
$SO(1,3)$, where $(1,3)$ is the signature of the Minkowski metric $(+,-,-,-)$.

Consider an infinitesimal Lorentz transformation
$$\Lambda\indices{^\mu_\nu} = \delta\indices{^\mu_\nu}+\omega\indices{^\mu_\nu},\quad
|\omega\indices{^\mu_\nu}|<<1\,\forall\mu,\nu
$$
where $\omega\indices{^\mu_\nu}$ are called \textit{parameters} of the transformation.

Substituting into the invariance condition yields
\begin{gather*}
g_{\mu\nu} \Lambda\indices{^\mu_\rho}\Lambda\indices{^\nu_\sigma} = g_{\rho\sigma}\\
g_{\mu\nu} \left(\delta\indices{^\mu_\rho}+\omega\indices{^\mu_\rho}\right)
\left(\delta\indices{^\nu_\sigma}+\omega\indices{^\nu_\sigma}\right)=g_{\rho\sigma}\\
g_{\mu\nu}\delta\indices{^\mu_\rho}\delta\indices{^\nu_\sigma} + 
\omega_{\nu\rho}\delta\indices{^\nu_\sigma} +
\omega_{\mu\sigma}\delta\indices{^\mu_\rho}=g_{\rho\sigma}\\
\omega_{\rho\sigma} = -\omega_{\rho\sigma}
\end{gather*}
so the parameters must be \textbf{antisymmetric}. This confirms the fact that
each transformation is uniquely determined by 6 quantities.
The transformation becomes
$$\Lambda\indices{^\mu_\nu}  =\delta\indices{^\mu_\nu} +
g^{\mu\rho}g\indices{^\sigma_\nu}\omega_{\rho\sigma}$$
since $\omega_{\rho\sigma}$ are antisymmetric, the only non-zero part of the product
$\omega_{\rho\sigma}g^{\mu\rho}g\indices{^\sigma_\nu}$ is given by the antisymmetric
part of $g^{\mu\rho}g\indices{^\sigma_\nu}$: anti-symmetrizing it yields
$$\Lambda\indices{^\mu_\nu} = \delta\indices{^\mu_\nu}+\frac{1}{2}\omega_{\rho\sigma}
\left(g^{\mu\rho}g\indices{^\sigma_\nu}-g^{\mu\sigma}g\indices{^\rho_\nu}\right)=
\delta\indices{^\mu_\nu} + i\frac{1}{2}\omega_{\rho\sigma}\left(M^{\rho\sigma}\right)
\indices{^\mu_\nu}
$$
confronting with the definition of generators
$$g=\mathbb{I} + i\alpha_aT^a$$
yields that the generators are
$$M^{\rho\sigma} = -i\left(g^{\mu\rho}g\indices{^\sigma_\nu}-g^{\mu\sigma}
g\indices{^\rho_\nu}\right)$$
so the generic Lorentz transformation can be expressed as
$$\Lambda = \exp\left[\frac{i}{2}\omega_{\rho\sigma}M^{\rho\sigma}\right]$$
and its commutation relations are
$$\left[M^{\lambda\tau},M^{\rho\sigma}\right]= -i \left(
    g^{\lambda\sigma}M^{\tau\rho}+g^{\tau\rho}M^{\lambda\sigma}-g^{\lambda\rho}
    M^{\tau\sigma}-g^{\tau\sigma}M^{\lambda\rho}
\right)$$
To make computations easier, the 6 generators are split into $3+3$ easier operators:
since $SO(3)$ is a subgroup of $SO(1,3)$, there is a way to express half of the
Lorentz Group generators in terms of those of $SO(3)$:
$$J_i  =\frac{1}{2}\varepsilon_{ijk}M^{jk}$$
which obey the $SO(3)$ algebra
$$\left[J_i,J_j\right] = i \varepsilon\indices{_i_j^k}J_k$$
and generate rotations; the boosts are generated by
$$K^i = M^{0i}$$
which obey the commutation relation
$$\left[K^i,K^j\right] = -i \varepsilon\indices{^i^j_k}J^k$$
which means that the algebra of the $K$'s is \textbf{not} closed, whereas the algebra
of the $J$'s is: this is due to the fact that 3D rotations form a subgroup of
$SO(1,3)$, while boosts do not (it is easy to check that the composition of 2 Cartesian
boosts is the composition of a Cartesian boost and a rotation).
The mixed commutator yields
$$[J_i,K_j] = i\varepsilon_{ijk}J^k$$
so the set of all $J$'s and all $K$'s is the set of the generators of $SO(1,3)$.
